\documentclass[12pt,a4paper]{article}
\usepackage[utf8]{inputenc}
\usepackage[T1]{fontenc}
\usepackage{lmodern}
\usepackage[french]{babel}
\usepackage{amsmath, amssymb, amsthm}
\usepackage{graphicx}
\usepackage{hyperref}
\usepackage{listings}
\usepackage{xcolor}
\usepackage{geometry}
\geometry{hmargin=2.5cm,vmargin=2.5cm}

\definecolor{codegreen}{rgb}{0,0.6,0}
\definecolor{codegray}{rgb}{0.5,0.5,0.5}
\definecolor{codepurple}{rgb}{0.58,0.0,0.82}
\definecolor{backcolour}{rgb}{0.95,0.95,0.92}

\lstdefinestyle{mystyle}{
    backgroundcolor=\color{backcolour},   
    commentstyle=\color{codegreen},
    keywordstyle=\color{magenta},
    numberstyle=\tiny\color{codegray},
    stringstyle=\color{codepurple},
    basicstyle=\ttfamily\footnotesize,
    breakatwhitespace=false,         
    breaklines=true,                 
    captionpos=b,                    
    keepspaces=true,                 
    numbers=left,                    
    numbersep=5pt,                  
    showspaces=false,                
    showstringspaces=false,
    showtabs=false,                  
    tabsize=2
}
\lstset{style=mystyle}

\title{\textbf{Rapport de Projet GPGPU : Méthodes de Monte-Carlo sur GPU (Version Allégée)}\\ \Large Master DSAD \& IA --- IUT de Douala}
\author{\textbf{Groupe :} [NOMS ET PRÉNOMS DES MEMBRES]}
\date{Date de remise : 10 Octobre 2026}

\begin{document}

\maketitle

\begin{abstract}
Ce projet présente l'implémentation et la parallélisation GPU d'algorithmes de Monte-Carlo pour l'estimation de \(\pi\), la valorisation d'options financières européennes sous le modèle de Black-Scholes et l'échantillonnage par chaîne de Markov (Metropolis-Hastings). La logique des noyaux a été développée et validée localement sur simulateur Numba-CUDA avant d'être portée sur GPU physique via Google Colab.
\end{abstract}

\section{Introduction et Contexte Mathématique}
Les méthodes de Monte-Carlo s'appuient sur la Loi des Grands Nombres pour approximer des espérances mathématiques par des moyennes empiriques.

\subsection{Convergence et Erreur Statistique}
Soit \(X_1, X_2, \dots, X_N\) des variables aléatoires indépendantes et identiquement distribuées. L'estimateur de Monte-Carlo est :
\begin{equation}
\bar{X}_N = \frac{1}{N} \sum_{i=1}^{N} X_i \xrightarrow[N \to \infty]{\text{p.s.}} \mathbb{E}[X]
\end{equation}
D'après le Théorème Central Limite, l'erreur statistique décroît selon \(\mathcal{O}\left(\frac{1}{\sqrt{N}}\right)\).

\subsection{Modèle de Black-Scholes}
Le prix d'un actif financier \(S_t\) suit un Mouvement Brownien Géométrique :
\begin{equation}
dS_t = r S_t dt + \sigma S_t dW_t
\end{equation}
La valeur exacte à maturité \(T\) pour un facteur d'incertitude \(Z \sim \mathcal{N}(0, 1)\) est :
\begin{equation}
S_T = S_0 \exp\left( \left(r - \frac{\sigma^2}{2}\right)T + \sigma\sqrt{T}Z \right)
\end{equation}
La formule fermée analytique d'un Call européen sert de référence pour la validation numérique :
\begin{equation}
C = S_0 \Phi(d_1) - K e^{-rT} \Phi(d_2)
\end{equation}
où \(d_1 = \frac{\ln(S_0/K) + (r + \sigma^2/2)T}{\sigma\sqrt{T}}\) et \(d_2 = d_1 - \sigma\sqrt{T}\).

\section{Architecture GPU Visée et Simulation Numba-CUDA}
L'architecture cible est basée sur le modèle SIMT (Single Instruction, Multiple Threads) des GPU NVIDIA. 

\subsection{Distribution de la Charge}
Chaque thread GPU exécute une boucle locale simulant \(N_{\text{thread}}\) trajectoires indépendantes en utilisant le générateur pseudo-aléatoire \texttt{xoroshiro128p\_normal\_float64}.
\begin{equation}
N_{\text{total}} = \text{Grille} \times \text{Blocs} \times N_{\text{thread}}
\end{equation}

\section{Validation Numérique et Résultats}
Les implémentations CPU (NumPy) et GPU (Numba-CUDA Simulé) ont été comparées.
\begin{table}[h!]
\centering
\begin{tabular}{|c|c|c|c|}
\hline
\textbf{Tirages (\(N\))} & \textbf{Prix MC CPU (€)} & \textbf{Prix MC GPU (€)} & \textbf{Écart / Exact (€)} \\ \hline
\(10^3\) & ... & ... & ... \\ \hline
\(10^5\) & ... & ... & ... \\ \hline
\(10^6\) & ... & ... & ... \\ \hline
\end{tabular}
\caption{Validation de la convergence vers la solution analytique de Black-Scholes.}
\end{table}

\section{Étape Finale : Exécution sur GPU Réel (Colab / Kaggle)}
\textit{[Présenter ici la mesure du speedup, la mémoire consommée et le nombre de simulations par seconde sur GPU T4/P100 réel.]}

\section{Conclusion}
Les méthodes de Monte-Carlo sont naturellement adaptées à la parallélisation massive sur GPU en raison de l'absence de dépendance mémoire entre les tirages.

\begin{thebibliography}{9}
\bibitem{black1973} Black, F., \& Scholes, M. (1973). The Pricing of Options and Corporate Liabilities. \textit{Journal of Political Economy}.
\bibitem{numba} Numba CUDA Documentation, \url{https://numba.readthedocs.io/}
\end{thebibliography}

\end{document}